Das Volumen einer Kugel ist proportional zur dritten Potenz des Durchmessers, $V \propto d^3$. Mit $V = V_0\left(1+\gamma\,\Delta\vartheta\right)$ wird
\begin{align}
 \frac{d}{d_0} &= \sqrt[3]{\frac{V}{V_0}} = \sqrt[3]{1+\gamma\,\Delta\vartheta}
\end{align}
Mit $\gamma = \gaHgO$ (Quecksilber, FoTaBe S.~191) und $\Delta\vartheta = \dTqF = \dTq$ ist die relative Änderung des Durchmessers
\begin{align}
 \frac{\Delta d}{d_0} &= \ddKF \\
   &= \sqrt[3]{1+\gaHg\cdot\dTq}-1 \\
   &= \ddK \approx \result{+\ddKP{0}{2}}
\end{align}
Für kleine $\gamma\,\Delta\vartheta$ gilt näherungsweise $\Delta d/d_0 \approx \frac{1}{3}\,\gamma\,\Delta\vartheta = \ddn$: Der Durchmesser ändert sich (wie jede Länge) relativ etwa ein Drittel so stark wie das Volumen.
